% GENERATED FILE. Edit canonical lessons, then run npm run generate:books.
% canonical-source-hash: fnv1a64:957dc95b70dc0279
% canonical-lessons: TE-C503-anuragam, TE-C503-mamakaram, TE-C503-apyayata, TE-C503-lalincu, TE-C503-abhimanam

\chapter{Love and Fondness}
\label{ch:te-love-and-fondness}
\hlchaptermodality{503}

\begin{chapteropening}
\textbf{By the end of this chapter:} \emph{I can name kinds of love and affection.}

Complete the last lesson of chapter 503: I can name kinds of love and affection.
\end{chapteropening}

\section[\texorpdfstring{\te{అనురాగం} (anurāgaṁ)}{అనురాగం (anurāgaṁ)}]{\te{అనురాగం} (anurāgaṁ) — tender love; affection}

\label{lesson:TE-C503-anuragam}



\begin{quote}
Before the new one: say the Telugu for to yearn, then the Telugu for to waste away.
\end{quote}

\subsection*{You'll want to know: \te{అనురాగం}}
\textbf{\te{అనురాగం}} — \emph{anurāgaṁ} — \textquotedblleft{}tender love; affection\textquotedblright{}.

\begin{quote}
\te{అక్కాచెల్లెళ్ళ} \te{మధ్య} \te{అనురాగం} \te{ఎక్కువ}.
\end{quote}

\begin{quote}
\emph{akkācelleḷḷa madhya anurāgaṁ ekkuva.}
\end{quote}

\begin{quote}
There is deep affection between the sisters.
\end{quote}

A warm, slightly literary word: \te{అనురాగంతో}, lovingly.

\begin{cousinweb}
\textbf{\te{అనురాగం}} (\emph{anurāgaṁ}) is borrowed from Sanskrit \emph{anurāga}, \textquotedblleft{}attachment, love\textquotedblright{} (CDIAL 331).
\end{cousinweb}

\begin{grammarlens}[title={What you've built}]
One more word for describing what you lived and felt.
\end{grammarlens}

\subsection*{Your turn}
\begin{itemize}
\raggedright
  \item \emph{Say it:} \emph{anurāgaṁ}
  \item \emph{Say it:} \emph{anurāgaṁ}, once more
  \item \emph{Say it:} \emph{akkācelleḷḷa madhya anurāgaṁ ekkuva.} — There is deep affection between the sisters.
  \item \emph{Recall:} say the Telugu for to yearn, then the Telugu for to waste away, then say \emph{anurāgaṁ} again
\end{itemize}

\begin{tcolorbox}[breakable,colback=teal!4,colframe=teal!35!black,title={Before you move on}]
What does \te{అనురాగం} mean? (\textquotedblleft{}Tender love; affection\textquotedblright{}.) Say it once more.
\end{tcolorbox}
\section[\texorpdfstring{\te{మమకారం} (mamakāraṁ)}{మమకారం (mamakāraṁ)}]{\te{మమకారం} (mamakāraṁ) — attachment; possessive love}

\label{lesson:TE-C503-mamakaram}



\begin{quote}
Before the new one: say the Telugu for to waste away, then the Telugu for tender love.
\end{quote}

\subsection*{You'll want to know: \te{మమకారం}}
\textbf{\te{మమకారం}} — \emph{mamakāraṁ} — \textquotedblleft{}attachment; possessive love\textquotedblright{}.

\begin{quote}
\te{ఆ} \te{పాత} \te{ఇంటి} \te{మీద} \te{నాన్నకి} \te{మమకారం} \te{ఎక్కువ}.
\end{quote}

\begin{quote}
\emph{ā pāta iṇṭi mīda nānnaki mamakāraṁ ekkuva.}
\end{quote}

\begin{quote}
Dad is very attached to that old house.
\end{quote}

What you are attached to takes \te{మీద}.

\begin{cousinweb}
\textbf{\te{మమకారం}} (\emph{mamakāraṁ}) is borrowed from Sanskrit \emph{mamakāra} (Monier-Williams), built from \emph{mama}, \textquotedblleft{}mine\textquotedblright{}, and \emph{kāra}, \textquotedblleft{}making\textquotedblright{}: the feeling that a person or a thing is your own.
\end{cousinweb}

\begin{grammarlens}[title={What you've built}]
One more word for describing what you lived and felt.
\end{grammarlens}

\subsection*{Your turn}
\begin{itemize}
\raggedright
  \item \emph{Say it:} \emph{mamakāraṁ}
  \item \emph{Say it:} \emph{mamakāraṁ}, once more
  \item \emph{Say it:} \emph{ā pāta iṇṭi mīda nānnaki mamakāraṁ ekkuva.} — Dad is very attached to that old house.
  \item \emph{Recall:} say the Telugu for to waste away, then the Telugu for tender love, then say \emph{mamakāraṁ} again
\end{itemize}

\begin{tcolorbox}[breakable,colback=teal!4,colframe=teal!35!black,title={Before you move on}]
What does \te{మమకారం} mean? (\textquotedblleft{}Attachment; possessive love\textquotedblright{}.) Say it once more.
\end{tcolorbox}
\section[\texorpdfstring{\te{ఆప్యాయత} (āpyāyata)}{ఆప్యాయత (āpyāyata)}]{\te{ఆప్యాయత} (āpyāyata) — warmth; loving care}

\label{lesson:TE-C503-apyayata}



\begin{quote}
Before the new one: say the Telugu for tender love, then the Telugu for attachment.
\end{quote}

\subsection*{You'll want to know: \te{ఆప్యాయత}}
\textbf{\te{ఆప్యాయత}} — \emph{āpyāyata} — \textquotedblleft{}warmth; loving care\textquotedblright{}.

\begin{quote}
\te{అమ్మమ్మ} \te{ఆప్యాయతతో} \te{అన్నం} \te{వడ్డించింది}.
\end{quote}

\begin{quote}
\emph{ammamma āpyāyatatō annaṁ vaḍḍiñcindi.}
\end{quote}

\begin{quote}
Grandma served the rice with loving care.
\end{quote}

\te{ఆప్యాయంగా}, affectionately, is the adverb.

\begin{cousinweb}
\textbf{\te{ఆప్యాయత}} (\emph{āpyāyata}) is usually traced to Sanskrit \emph{āpyāyana}, \textquotedblleft{}causing fullness; gladdening\textquotedblright{} (Monier-Williams), with the abstract ending -\te{త}: the warmth that makes others feel cared for.
\end{cousinweb}

\begin{grammarlens}[title={What you've built}]
One more word for describing what you lived and felt.
\end{grammarlens}

\subsection*{Your turn}
\begin{itemize}
\raggedright
  \item \emph{Say it:} \emph{āpyāyata}
  \item \emph{Say it:} \emph{āpyāyata}, once more
  \item \emph{Say it:} \emph{ammamma āpyāyatatō annaṁ vaḍḍiñcindi.} — Grandma served the rice with loving care.
  \item \emph{Recall:} say the Telugu for tender love, then the Telugu for attachment, then say \emph{āpyāyata} again
\end{itemize}

\begin{tcolorbox}[breakable,colback=teal!4,colframe=teal!35!black,title={Before you move on}]
What does \te{ఆప్యాయత} mean? (\textquotedblleft{}Warmth; loving care\textquotedblright{}.) Say it once more.
\end{tcolorbox}
\section[\texorpdfstring{\te{లాలించు} (lāliñcu)}{లాలించు (lāliñcu)}]{\te{లాలించు} (lāliñcu) — to fondle; to cuddle}

\label{lesson:TE-C503-lalincu}



\begin{quote}
Before the new one: say the Telugu for attachment, then the Telugu for warmth.
\end{quote}

\subsection*{You'll want to know: \te{లాలించు}}
\textbf{\te{లాలించు}} — \emph{lāliñcu} — \textquotedblleft{}to fondle; to cuddle\textquotedblright{}.

\begin{quote}
\te{ఏడుస్తున్న} \te{బాబును} \te{అమ్మ} \te{లాలించింది}.
\end{quote}

\begin{quote}
\emph{ēḍustunna bābunu amma lāliñcindi.}
\end{quote}

\begin{quote}
Mum cuddled the crying baby boy.
\end{quote}

Past (I …): \te{లాలించాను}, \emph{lāliñcānu}.

Also to soothe and humour a child.

\begin{cousinweb}
\textbf{\te{లాలించు}} (\emph{lāliñcu}) is built on Sanskrit \emph{lālana}, \textquotedblleft{}caressing\textquotedblright{} (CDIAL 11025), with the Telugu verb ending -\te{ఇంచు}.
\end{cousinweb}

\begin{grammarlens}[title={What you've built}]
One more word for describing what you lived and felt.
\end{grammarlens}

\subsection*{Your turn}
\begin{itemize}
\raggedright
  \item \emph{Say it:} \emph{lāliñcu}
  \item \emph{Say it:} \emph{lāliñcu}, once more
  \item \emph{Say it:} \emph{ēḍustunna bābunu amma lāliñcindi.} — Mum cuddled the crying baby boy.
  \item \emph{Recall:} say the Telugu for attachment, then the Telugu for warmth, then say \emph{lāliñcu} again
\end{itemize}

\begin{tcolorbox}[breakable,colback=teal!4,colframe=teal!35!black,title={Before you move on}]
What does \te{లాలించు} mean? (\textquotedblleft{}To fondle; to cuddle\textquotedblright{}.) Say it once more.
\end{tcolorbox}
\section[\texorpdfstring{\te{అభిమానం} (abhimānaṁ)}{అభిమానం (abhimānaṁ)}]{\te{అభిమానం} (abhimānaṁ) — fondness; regard}

\label{lesson:TE-C503-abhimanam}



\begin{quote}
Before the new one: say the Telugu for warmth, then the Telugu for to fondle.
\end{quote}

\subsection*{You'll want to know: \te{అభిమానం}}
\textbf{\te{అభిమానం}} — \emph{abhimānaṁ} — \textquotedblleft{}fondness; regard\textquotedblright{}.

\begin{quote}
\te{ఆ} \te{నటుడంటే} \te{నాకు} \te{చాలా} \te{అభిమానం}.
\end{quote}

\begin{quote}
\emph{ā naṭuḍaṇṭē nāku cālā abhimānaṁ.}
\end{quote}

\begin{quote}
I am very fond of that actor.
\end{quote}

Also self-respect or pride, depending on the sentence.

\begin{cousinweb}
\textbf{\te{అభిమానం}} (\emph{abhimānaṁ}) is borrowed from Sanskrit \emph{abhimāna} (CDIAL 524), which meant pride and self-regard; in Telugu it is also the warm regard you feel for a person.
\end{cousinweb}

\begin{grammarlens}[title={What you've built}]
One more word for describing what you lived and felt.
\end{grammarlens}

\subsection*{Your turn}
\begin{itemize}
\raggedright
  \item \emph{Say it:} \emph{abhimānaṁ}
  \item \emph{Say it:} \emph{abhimānaṁ}, once more
  \item \emph{Say it:} \emph{ā naṭuḍaṇṭē nāku cālā abhimānaṁ.} — I am very fond of that actor.
  \item \emph{Recall:} say the Telugu for warmth, then the Telugu for to fondle, then say \emph{abhimānaṁ} again
\end{itemize}

\begin{tcolorbox}[breakable,colback=teal!4,colframe=teal!35!black,title={Before you move on}]
What does \te{అభిమానం} mean? (\textquotedblleft{}Fondness; regard\textquotedblright{}.) Say it once more.
\end{tcolorbox}
